\section{はじめに}

近年普及が加速しているIoTデバイスの多くは\cite{makhshari2021iot}，REST (Representational State Transfer) アーキテクチャスタイル\cite{fielding2000architectural}に基づいたREST APIを提供している．クライアント開発者はAPI仕様を満たす必要があるが，レスポンスのエラーメッセージは情報が少なく発生箇所を示さないため，繰り返し送受信をしなければならない．
先行研究\cite{inoue2025dbrag}において，REST API仕様をベクトルデータベースへ格納し，RAG (Retrieval Augmented Generation)\cite{lewis2020retrieval}で関連仕様を検索してLLMへ与える手法を検証した．非推奨仕様と最新仕様をそれぞれコード片と自然言語テキストに分けて4つのデータベースへ格納する構成が最も高い性能を示した．一方で，2つの課題が残っている．
第一に，仕様を分割して格納することで，非推奨仕様と最新仕様の対応関係が検索結果に直接反映されず，両者の対応付けをLLMの推論に頼る必要がある．第二に，データベースを分割するほど検索結果の総量が増えるため，構造化による効果とコンテキスト量による効果を切り分けるのが難しい．

そこで本研究では，検索単位を仕様のテキスト片ではなく仕様の差分に置く手法を提案する．非推奨仕様の要素から最新仕様の要素への対応をエッジとして知識グラフで表現し，誤用コードから抽出した要素で検索する．
さらに対応付け先から必要な要素を多ホップ探索することで，修正に必要な要素群をまとめて取得する．


% これ消すか，マージするか検討！！
REST APIに関する知識をグラフとして表現する試みとして，API制約の知識グラフによる誤用検出\cite{ren2020api}や，文書のコミュニティ要約を検索単位とするGraph RAG\cite{edge2024graphrag}がある．これらはいずれも単一時点の知識を対象としており，バージョン間の対応関係を検索単位とはしていない．本研究では仕様間の差分を検索単位とする．


% 本研究では以下のRQを設定した．
% \begin{description}[leftmargin=1.5em, itemsep=0pt, topsep=2pt]
%   \item[RQ1:] {\bf 仕様の差分をグラフ構造で与えることは，仕様本文をそのまま与える場合より修正率を向上させるか？}
%   \item[RQ2:] {\bf 提案手法のどの構成要素が修正率に寄与しているか？}
% \end{description}
